\cleardoublepage
\chapternonum{致谢}

时间如白驹过隙，在MCLab、浙江大学生命科学学院和“农学+X”多学科交叉人才培养卓越中心的支持下我稳步前行。如今，二十一年的求学生涯即将迎来尾声，我也即将褪去学生的身份前往更加广阔的天地。值此论文完成之际，无数回忆涌上心头，我想感谢一路走来给予过我关怀和帮助的人。

特别感谢我的导师陈铭教授，陈老师对我的支持和信任是我能够顺利完成博士学业的关键。我的工位面朝陈老师办公室，因此，总能看到陈老师微笑地关心我的科研和生活。我们讨论课题、申奖立项、规划未来；也曾在科研之余亲近自然、放松身心。我始终难忘与陈老师共同经历的许多宝贵时光：三次前往新西伯利亚开展学术交流，穿越丛林欣赏鄂毕河的美景；在圣彼得堡感受历史与人文；纵览亚欧交接的伊斯坦布尔；感受古巴的热带风光和内蒙古呼伦贝尔大草原的辽阔壮美。这些经历不仅开阔了我的学术视野，也让我对科研、人生与世界有了更深层次的理解。更令我受益终身的，是陈老师严谨求实的学术态度、卓越的学术洞察力以及对学科发展的深刻思考。无论何时，陈老师都充分尊重我的科研兴趣，并总能在关键时刻指出问题核心，给予清晰而富有启发性的指导。此外，我也特别感谢陈老师的宽容、理解与鼓励，使我能够在南京大学三年的学习与科研过程中始终坚定前行。陈老师的言传身教与谆谆教诲，将始终激励我不断前进。同时，衷心感谢李燕老师在学习和生活中给予我的关心与帮助，感谢李老师长期以来为实验室所付出的辛勤努力与无私奉献。特别是在我南京交流期间，李老师在生活上给予了我许多照顾与关怀，使我深受感动，在此致以诚挚的感谢。

感谢陈迪俊教授。迪俊师兄始终是我科研道路上的楷模，他勤勉踏实、严谨专注、待人真诚，令我深感敬佩。在科研过程中，我曾多次见到迪俊师兄夜以继日地工作，面对困难和挑战从不退缩。与此同时，迪俊师兄扎实的学术功底、敏锐的问题意识以及出色的科研能力，使我受益良多。能够在科研与学习过程中得到迪俊师兄的帮助与指导，我深感幸运与感激。同样，我也会永远牢记在南京生活的三年时光，这是一段与我自身而言，人生观与价值观成长最为迅速、改变最大、记忆最为难忘的旅程。此外，也特别感谢傅靓彧老师在学习与生活中给予我的关心与支持，让我感受到许多温暖与鼓励。

感谢李卓瑾，她作为我的女朋友给予了我莫大的关怀与陪伴。她知识面广、单纯善良、默默付出。她是我面对周遭生活的底气，使我的内心更加坚定，不惧世俗。感谢她对我科研上的督促、鼓励与支持，我们一起经历了许多美好的时光，分享了许多快乐和困难。她的存在让我感到幸福和满足，我也非常珍惜我们之间的感情。同样感谢我们的猫咪悠悠，悠悠性格稳定，活泼可爱，给我们的生活带来了很多乐趣和温暖。

感谢冯聪师兄，他是MCLab的大师兄、顶梁柱，也是我的良师益友，是我今生今世永远都会牢记于心、永远不会遗忘的师兄。感谢王晶晶和张霈婧师姐、陈宏俊、胡言石、李思达和陈俞皓师兄对我科研道路上的指点迷津，他们像一本鲜活的百科全书，总能给我带来新的启发和思路。他们也会把MCLab经营的很好，张罗包粽子、捏饺子、过生日，让我感受到实验室的温暖和团结。感谢胡跃明、倪清扬、赵亮、辛赛格同门的一路陪伴。跃明是在浙大陪伴我时间最久、最亲切的人。他阅历丰富，看待问题鞭辟入里，与他同行倍感荣幸。清扬是我在浙大同门同级中非常亲密的朋友，我们一起经历了研究生生涯的酸甜苦辣。同时，感谢张世龙、王子轩和吴梁炜师弟、刘恩言、刘丽亚、陈一凡和谢璐瑶师妹给实验室带来欢快的氛围，他们是未来的希望。

感谢南京大学的阮忠豪师弟，他乐观开朗、潇洒能干，是我科研道路上不可多得的伙伴。感谢于冉冉师姐，她心思细腻，永远会在饭卡中留钱让我使用，我会永远牢记在心。感谢蓝杨铭师姐，她细心体贴、开朗大方、处事淡然，有她在的地方总是充满欢声笑语，其乐融融。感谢曹广硕师兄，他为人坦荡，内敛但热心。感谢吕同轩，他性情爽朗，为人正直，能力出众。只要有同轩在的地方就会有聊不完的天。感谢唐一骏师弟，他是位妙趣横生、金句频出的师弟，也有着非常强的责任心和执行力。总之和他们谈笑风生保证了我的身心健康。同时，也感谢赵雪师姐和张晨露师妹对我科研上的协助。

最后，我想感谢我的父亲晁道松，他勤劳善良、正直诚实、动手能力超强、喜欢唱歌，带领家庭从农村走出来，并给予了我无尽的爱和支持。感谢我的母亲段利萍，她勤劳能干、乐观开朗，爱孩子胜过爱自己，现在的她退休后终于能好好享受自己的人生了。感谢我的姐姐晁好月，她是陪我一起长大的月月姐，我们打打闹闹，共同成长，是她在前面探路才有了我顺利的求学生涯。总之，他们为我构建了求学之路上最坚实的后盾和保障。未来的路还有很长，希望未来自己在科研上能有所建树，成就自己也回馈家人。

\hfill 晁好瑜

\hfill 二零二六年五月


